\documentclass[notitlepage,a4paper]{article}
\usepackage[top=3cm,bottom=3cm,left=2cm,right=2cm]{geometry} %paper/margins


\usepackage[bitstream-charter]{mathdesign} 
%\usepackage{charter}
%\usepackage{beramono}


\usepackage[parfill]{parskip} 

%\numberwithin{equation}{section}

\usepackage{enumerate}
\usepackage{lipsum}


%\usepackage[explicit]{titlesec}													%Question 1
%\titleformat{\section}{\normalfont\Large\bfseries}{}{0em}{#1\ \thesection}

%----------------------------------------------------------------------------------------
%	href/url
%----------------------------------------------------------------------------------------

\usepackage{hyperref}
\hypersetup{
    colorlinks=true,
    linkcolor=blue,
    filecolor=magenta,      
    urlcolor=blue,
    pdftitle={Overleaf Example},
    pdfpagemode=FullScreen,
    }

%----------------------------------------------------------------------------------------
%	date format
%----------------------------------------------------------------------------------------

\usepackage[yyyymmdd]{datetime}
\renewcommand{\dateseparator}{--}

%----------------------------------------------------------------------------------------
%	page header
%----------------------------------------------------------------------------------------

\usepackage{fancyhdr}
 
\pagestyle{fancy}
\fancyhf{}
\lhead{SURE WP3 - Energy Model}
\rhead{Toby Simpson \today}
\cfoot{\thepage}

%----------------------------------------------------------------------------------------
%	graphics
%----------------------------------------------------------------------------------------

\usepackage{graphicx}
\graphicspath{ {img/} }

\usepackage{placeins} 			%\FloatBarrier


%\begin{figure}[h]
%\centering
%\includegraphics[trim={110 0 110 0},clip,width=\textwidth]{fig1}
%\caption{Richardson without preconditioner}
%\end{figure}


\usepackage{tikz}
\usetikzlibrary{shapes,arrows,fit}
\usetikzlibrary{arrows.meta}
\usetikzlibrary{patterns}


%----------------------------------------------------------------------------------------
%	math shortcuts
%----------------------------------------------------------------------------------------

\newcommand{\x}{\mathbf{x}} 
\renewcommand{\u}{\mathbf{u}}    
\renewcommand{\v}{\mathbf{v}}   
\renewcommand{\b}{\mathbf{b}}      
\newcommand{\f}{\mathbf{f}}      
\newcommand{\n}{\mathbf{n}}     
\renewcommand{\r}{\mathbf{r}}     

\newcommand{\A}{\mathbf{A}}    
\newcommand{\E}{\mathbf{E}}   
\newcommand{\F}{\mathbf{F}}   
\newcommand{\X}{\mathbf{X}}   
\renewcommand{\P}{\mathbf{P}}   
\renewcommand{\S}{\mathbf{S}}   


\newcommand{\R}{\mathbb{R}}
\newcommand{\N}{\mathbf{N}}
\newcommand{\I}{\mathbf{I}}          

\renewcommand{\hat}{\widehat}   
\renewcommand{\tilde}{\widetilde}   


%----------------------------------------------------------------------------------------
%	document
%----------------------------------------------------------------------------------------


\begin{document}


\tableofcontents
\newpage

%===============================
\section{Introduction}
%=============================== 

%===============================
\subsection{Overview}
%=============================== 

This document describes the function and use of the website: \url{http://toby.euler.usi.ch}.

The website allows users to generate scenarios for energy transition in Switzerland and to measure and compare them according to set of metrics. The decisions of the users are stored and can be further analysed by researchers.

The model provides the underlying calculations and Application Programmer Interface (API) for the SURE WP3 game, but can also be used in a stand alone manner. Scenarios generated by the model are stored for comparison and analysis which forms part of SURE WP11. 

The site does not currently require user authentication and its use should be limited to trusted users.

%===============================
\subsection{Representative Energy System}
%=============================== 

The Representative Energy System (RES) provides a simple model the flow of energy over time through the Swiss economy. The structure is summarasied in Figure \ref{fig:res1}. Energy flows through three basic stages, which are outlined below.


\begin{figure}[ht]

\tikzstyle{txt1} = [rotate=90, anchor=west]
\tikzstyle{arr1} = [thick,-Latex]
\tikzstyle{arr2} = [thick, -{Circle[scale=0.5]}, shorten >=-0.4mm]
\tikzstyle{arr3} = [thick, {Circle[scale=0.5]}-Latex, shorten <=-0.4mm]
\tikzstyle{c1}	 = [blue]
\tikzstyle{c2}	 = [violet]
\tikzstyle{c3}	 = [magenta]
\tikzstyle{c4}	 = [cyan]
\tikzstyle{c5}	 = [teal]

\tikzstyle{input}   = [rectangle, draw, text width=3cm, minimum height=0.75cm]
\tikzstyle{demand}  = [rectangle, draw, text width=3cm, minimum height=0.6cm]


\begin{tikzpicture}[scale=0.45, y=-1cm]

%	\draw[help lines] (0,0) grid (37,36);
	
	%imports
	\draw (04,01) node[anchor=south] {\texttt{Imported}};
	\draw (00,01) rectangle (09,14) ;
	
	\draw (01,02) rectangle (08,04) node[pos=0.5] (ele) {Electricity};
	\draw (01,05) rectangle (08,07) node[pos=0.5] (pet) {Petroleum};
	\draw (01,08) rectangle (08,10) node[pos=0.5] (gas) {Natural Gas};
	\draw (01,11) rectangle (08,13) node[pos=0.5] (nuc) {Nuclear Fuel};
	
	%domestic
	\draw (04,17) node[anchor=south] {\texttt{Domestic}};
	\draw (00,17) rectangle (09,33);
	
	\draw (01,18) rectangle (08,20) node[pos=0.5] (res) {River Flow};
	\draw (01,21) rectangle (08,23) node[pos=0.5] (riv) {Reservoir Storage};
	\draw (01,24) rectangle (08,26) node[pos=0.5] (sol) {Solar Energy};
	\draw (01,27) rectangle (08,29) node[pos=0.5] (wnd) {Wind Energy};
	\draw (01,30) rectangle (08,32) node[pos=0.5] (bio) {Biomass/Waste};
	
	%demand
	\draw (32,08) node[anchor=south] {\texttt{Final Demand}};
	\draw (28,08) rectangle (37,33);

	%demand
	\draw (29,09) rectangle (36,12) node[pos=0.5] {Base Electricity};
	\draw (29,13) rectangle (36,16) node[pos=0.5] {Cooling Electricity};
	\draw (29,17) rectangle (36,20) node[pos=0.5] {Industrial Heat};
	\draw (29,21) rectangle (36,24) node[pos=0.5] {Residential Heat};
	\draw (29,25) rectangle (36,28) node[pos=0.5] {Road Transport};
	\draw (29,29) rectangle (36,32) node[pos=0.5] {Rail Transport};
	
	%htrans
	\draw[arr2,c1] (08,03) -- (21,03);
	\draw[arr2,c2] (08,06) -- (22,06);
	\draw[arr2,c3] (08,09) -- (23,09);
	\draw[arr2,c4] (08,12) -- (13,12);
	
	%vgen
	\draw[arr3,c3] (15,09) -- (15,18);
	\draw[arr3,c4] (13,12) -- (13,18);
	
	%gen
	\draw[pattern=north east lines, pattern color=gray!50] (12,18) rectangle (16,32) node[pos=0.5, rotate=90, fill=white] {Electricity Generation};
	
	%gen
	\draw[arr1,c5] (08,19) -- (12,19);
	\draw[arr1,c5] (08,22) -- (12,22);
	\draw[arr1,c5] (08,25) -- (12,25);
	\draw[arr1,c5] (08,28) -- (12,28);
	\draw[arr1,c5] (08,31) -- (12,31);
	
	\draw[arr2,c1] (16,24) -- (21,24);
	
	%distribution
	\draw (22,01) node[anchor=south] {\texttt{Distribution}};
	\draw (19,01) rectangle (25,33);
	
	%vdist
	\draw[thick,c1] (21,03) -- (21,30);
	\draw[thick,c2] (22,06) -- (22,31);
	\draw[thick,c3] (23,09) -- (23,23);
	
	%con
	\draw[arr3,c1] (21,10.5) -- (29,10.5);
	\draw[arr3,c1] (21,14.5) -- (29,14.5);
	
	\draw[arr3,c1] (21,18) -- (29,18);
	\draw[arr3,c3] (23,19) -- (29,19);
	
	\draw[arr3,c1] (21,22) -- (29,22);
	\draw[arr3,c3] (23,23) -- (29,23);
	
	\draw[arr3,c1] (21,26) -- (29,26);
	\draw[arr3,c2] (22,27) -- (29,27);
	
	\draw[arr3,c1] (21,30) -- (29,30);
	\draw[arr3,c2] (22,31) -- (29,31);
	

\end{tikzpicture}
\caption{Representative Energy System}
\label{fig:res1}
\end{figure}



%===============================
\subsubsection{Primary Energy}
%===============================

The inputs of energy for the model are either imported or domestically available.  Imported energy is as follows, including a three letter code which is used in parameter naming:

\begin{description}
\item[imp] Imported electricity
\item[pet] Petroleum products, including petrol, diesel and oil
\item[gas] Natural gas
\item[nuc] Nuclear fuel
\end{description}

Domestic sources of energy along with their codes are as follows:

\begin{description}
\item[riv] River flow, the seasonally varying source of hydroelectric power 
\item[res] Reservoir storage, the non-seasonal source of hydroelectric power
\item[sol] Solar energy, seasonally varying source for photovoltaic electricity generation
\item[wnd] Wind energy, seasonally available source for wind turbines
\item[bio] Biomass/waste the non-seasonal energy available for combustion in the form of bio fuels such as wood or waste
\end{description}


%===============================
\subsubsection{Transformation}
%===============================

This includes the conversion and delivery of energy in various form in order to satisfy demand. There are two main types:

\begin{itemize}
\item Generation of electricity from the primary inputs, requiring infrastructure such as reservoirs, solar panels, wind turbines and combustion plants.
\item Distribution of fuel to end-users, such as petroleum for road transport and gas for use in domestic and industrial heating, again requiring infrastructure such as pipelines, and power transmission lines.
\end{itemize}

%===============================
\subsubsection{Final Demand}
\label{sec:final}
%===============================

This represents the consumption of energy by end-users as defined below.  The naming convention of the model gives each a four letter code which is used in parameter names where relevant:

\begin{description}
\item[base] Base demand is the non-seasonal demand for electricity for both residential and business use, such as lighting, appliances and communications.
\item[cool] Cooling is the seasonally dependent demand for electricity for use in refrigeration and air conditioning for residential and business use.
\item[hind] Industrial Heat is the seasonally-independent demand for either natural gas or electricity to provide heat for industrial processes such as manufacturing.
\item[hres] Residential heat is the seasonally-dependent demand for either natural gas or electricity for domestic space and water heating.  
\item[road] Road Transport is the non-seasonal use of petroleum products or electricity for transportation by road.
\item[rail] Rail Transport is the non-seasonal use of petroleum products or electricity for transportation by rail.
\end{description}


%===============================
\subsubsection{Costs \& Externalities}
%===============================

At each stage of system there may be requirements for for fuel and infrastructure which incur fixed and floating costs, as well as externalities such as emissions, land use and the production of waste. The user can allocate capacity to the various different energy pathways and the model calculates the resulting values for cost and externalities according to pre-defined weights.

%===============================
\subsubsection{Seasonality \& Substitution}
%===============================

The supply of energy inputs may not be constant, since resources such as river flow, solar and wind energy vary seasonally. Similarly some elements of demand will fluctuate seasonally, such as residential heating and cooling electricity. 

As can be seen in Figure \ref{fig:res1}, some elements of final demand may be satisfied either by electricity of fossil fuels.  In these cases, the user can define a values for the substitution of one type of consumption for another.  For example as consumers buy and use more electric cars the demand for road transport will increasingly be satisfied via electricity rather that petroleum.

%===============================
\subsubsection{Time, Metrics \& Analysis}
%===============================

The model considers a period of time over which the simulation takes place, and an interval at which the model is re-calculated.  The model considers intrinsic and user-defined fluctuations in supply and demand, and calculates averaged values for cost and externalities on an annualised basis. The averaged values can be used to compare the outcomes of different instances of the simulation. 

Finally, since all of the decisions made by the users are recorded per simulation, they can be made available for further analysis by researchers who can quantify their decisions with regard to the various technological, financial and environmental factors involved.


%===============================
\section{Definitions \& Calculation}
%=============================== 

The software environment for the model makes use of a database and web server on a standard LAMP architecture.  Users interact via a set of php web pages which submit data to an SQL database which then makes calculations and returns results to the browser.  In addition, as part of SURE WP3 there is an API comprising of a set of HTTP calls by which allow programmatic interaction with the model, with responses returned as XML documents.

Fundamental to the model is the database, which contains a set of parameters and functions.  The parameters are listed in full in Appendix \ref{app:prm1}. Their values may be pre-defined by the model designers, user-defined per simulation or calculated according to the definitions of the model.  

Parameters are grouped into categories which reflect the general structure of the model. Both the website itself and following sections use the same categorical structure to define the parameters and their interactions.

%===============================
\subsection{Time}
\label{sec:time}
%=============================== 

The model performs integration over a time period which can be defined per RES in the settings section. There are three parameters:

\begin{description}
\item[res\_dt] is the length of the time step in years.  The default value is 0.0833333, which is $\frac{1}{12}$  of a year or one notional month.
\item[res\_nt] is the total number of time steps in the calculation. These will be referred to as periods and are numbered from zero. The total time in years for a simulation is therefore given as \texttt{res\_dt} * \texttt{res\_nt}. Since the default value is 120 the default simulation time is 10 years.
\item[res\_mod] is the interval at which a parameter is available for editing. When a parameter is not calculated the user can set its value at any time period. For simplicity the user interface will only show edit controls at every \texttt{res\_mod} period.  Since the default value is 6 the user interface will allow editing at every: \texttt{res\_mod} * \texttt{res\_dt} = 0.5 years, or six notional months.
\end{description}

At any period during the simulation the value for time \texttt{t} in years is given by the cumulative sum of \texttt{res\_dt} or the product of the period number and \texttt{res\_dt} (which yield the same result).

%===============================
\subsection{Environment}
%=============================== 

%===============================
\subsubsection{Society}
\label{sec:pop}
%=============================== 

This section contains parameters and a calculated value for the population. They are defined as follows:

\begin{description}
\item[soc\_pop\_ini] is the initial size of the population in millions.
\item[soc\_pop\_grw] is the population growth rate factor per year. This leads to exponential growth as defined below.
\item[soc\_pop\_mig] is the net inward migration rate per year in millions.
\item[soc\_pop] is the population size in millions.  The population growth over a period $n$ is given as the integration of an ordinary differential equation (ODE) as follows:
\begin{eqnarray}
\texttt{soc\_pop}_{n+1} = \texttt{soc\_pop}_n + \texttt{res\_dt} * (\texttt{soc\_pop}_n * \texttt{soc\_pop\_grw}_n + \texttt{soc\_pop\_mig}_n)
\end{eqnarray}
\end{description}

Other economic and demographic factors may be added in future, specifically those relating to economic growth and energy efficiency.

%===============================
\subsection{Demand}
%===============================

%===============================
\subsubsection{Demand per Capita \& Final Demand}
%=============================== 

The size of the population of Section \ref{sec:pop} determines the quantities of the various forms of final demand as defined in Section \ref{sec:final}. The user-defined parameters in the demand per capita category (or requirement) are multiplied by the population size to give values for type each of final demand.  For example:
\begin{eqnarray}
\texttt{dem\_base} &=& \texttt{soc\_pop} * \texttt{req\_base} \\
\texttt{dem\_cool} &=& \texttt{soc\_pop} * \texttt{req\_cool} \\
				&\vdots& \\
\texttt{dem\_rail} &=& \texttt{soc\_pop} * \texttt{req\_rail} 
\end{eqnarray}

The units of final demand are given as annualised units of generated electricity (GWh). In general this flows through to all subsequent calculations unless otherwise indicated.
 
 
%===============================
\subsubsection{Electrification}
%=============================== 

Final demand is satisfied by technology that consumes energy in different forms. The process of electrification reflects the substitution of technology that consumes fossil fuels such as petroleum products and natural gas by technology that consumes electricity.  An example would be the purchase of electric cars over petrol ones in the satisfaction of road transport demand. The process of electrification applies to the four categories of demand (industrial heat, residential heat, road transport and rail transport) which can consume both fossil fuels and electricity.

Electrification is calculated as a scalar value between zero and one, where zero indicates all fossil fuel use and one indicates all electricity. Since the process of the substitution of technology takes place constantly over time it is modelled as a linear trend with an initial value and an annual rate of change. For example, the user defined electrification parameters for road transport are as follows:

\begin{description}
\item[sub\_road\_ini] is the initial value for electrification.  The default is 0.2 which indicates that 20\% of demand for road transport is satisfied by electric vehicles and the remaining 80\% comes from vehicles powered by fossil fuels.
\item[sub\_road\_lin] is the linear trend in electrification. The default value of 0.01 indicates that each year 1\% of vehicles that consume fossil fuel are replaced by electric vehicles. 
\item[sub\_road] is the value of electrification calculated by the linear extrapolation as a function of time in years as defined in Section \ref{sec:time}, for example with road transport:
\begin{eqnarray}
\texttt{sub\_road} &=& \texttt{sub\_road\_ini} + (\texttt{t} * \texttt{sub\_road\_lin}) 
\end{eqnarray}
\end{description}


%===============================
\subsubsection{Consumption}
%=============================== 

Given values for final demand and electrification, the consumption of each type of demand and type of energy can be calculated.  For electricity-only demand consuption is equal to final demand:
\begin{eqnarray}
\texttt{con\_base\_ele} &=& \texttt{dem\_base} \\
\texttt{con\_cool\_ele} &=& \texttt{dem\_cool}
\end{eqnarray}
Electricity consumption for other types of demand is calculated using substitution:
\begin{eqnarray}
\texttt{con\_hind\_ele} &=& \texttt{dem\_hind} *  \texttt{sub\_hind}\\
\texttt{con\_hres\_ele} &=& \texttt{dem\_hres} *  \texttt{sub\_hres}\\
\texttt{con\_road\_ele} &=& \texttt{dem\_road} *  \texttt{sub\_road}\\
\texttt{con\_rail\_ele} &=& \texttt{dem\_rail} *  \texttt{sub\_rail}
\end{eqnarray}
Finally, the fossil fuel demand is calculated as the remainder that is not substituted:
\begin{eqnarray}
\texttt{con\_hind\_gas} &=& \texttt{dem\_hind} *  (1 - \texttt{sub\_hind})\\
\texttt{con\_hres\_gas} &=& \texttt{dem\_hres} *  (1 - \texttt{sub\_hres})\\
\texttt{con\_road\_pet} &=& \texttt{dem\_road} *  (1 - \texttt{sub\_road})\\
\texttt{con\_rail\_pet} &=& \texttt{dem\_rail} *  (1 - \texttt{sub\_rail})
\end{eqnarray}

As a simplification, it is assumed that all fossil demand for industrial and residential heat comes from gas, and all fossil demand for transport is petroleum.  It is important to note that substitution of one technology for another may lead to significant increases in efficiency.  For instance it is estimated that electric cars consume three times less energy than theor petrol equivalents.  Tis is an important omission from the model which will be addressed in future.  

Finally, the different types of consumption can be aggregated by energy type to give consumption totals as follows:
\begin{eqnarray}
\texttt{con\_ele} &=& (\texttt{con\_base\_ele} + \texttt{con\_cool\_ele} \\
&+& \texttt{con\_hind\_ele} + \texttt{con\_hres\_ele} \\
&+& \texttt{con\_road\_ele} + \texttt{con\_rail\_ele} )\\
\texttt{con\_gas} &=& \texttt{con\_hind\_gas} + \texttt{con\_hres\_gas}\\
\texttt{con\_pet} &=& \texttt{con\_road\_pet} + \texttt{con\_rail\_pet}
\end{eqnarray}

It should be clear from above that as time progresses, increasing electrification will decrease demand for fossil fuels and increase demand for electricity.

%===============================
\subsection{Generation}
%=============================== 

%===============================
\subsection{Imports}
%=============================== 

%===============================
\subsection{Metrics}
%=============================== 


\appendix

%===============================
\section{Appendix}
%=============================== 


%===============================
\subsection{Settings}
%=============================== 

%===============================
\subsection{Parameters}
\label{app:prm1}
%=============================== 

List of parameters

\begin{table}[h!]
\begin{center}
\begin{tabular}{ccc}
 cell1 & cell2 & cell3 \\ 
 \hline
 cell4 & cell5 & cell6 \\  
 cell7 & cell8 & cell9    
\end{tabular}
\end{center}
\end{table}



%===============================
\subsection{Visualisation/Data}
%=============================== 


%===============================
\subsection{API}
%=============================== 




\end{document}